\documentclass[11pt]{article}
\usepackage[T1]{fontenc}
\usepackage{lmodern}
\usepackage[a4paper,margin=2.2cm]{geometry}
\usepackage{xcolor}
\usepackage{listings}
\usepackage[skins]{tcolorbox}
\usepackage{titlesec}
\usepackage{enumitem}
\usepackage{booktabs}
\usepackage{fancyhdr}
\usepackage[hidelinks]{hyperref}

% Assignment details.
\newcommand{\course}{CS 150 Data Structures}
\newcommand{\assignment}{Assignment 2: Hash Tables}
\newcommand{\student}{Kai Turner (ID 20260417)}

\definecolor{codebg}{RGB}{248,248,250}
\definecolor{codekw}{RGB}{0,90,170}
\definecolor{codestr}{RGB}{160,60,20}
\definecolor{codecom}{RGB}{100,120,100}
\definecolor{head}{RGB}{40,40,60}

% Code style: change language or colours here.
\lstdefinestyle{assignment}{
  language=Python,
  basicstyle=\ttfamily\small,
  keywordstyle=\color{codekw}\bfseries,
  stringstyle=\color{codestr},
  commentstyle=\color{codecom}\itshape,
  backgroundcolor=\color{codebg},
  numbers=left, numberstyle=\tiny\color{gray}, numbersep=8pt,
  frame=leftline, framerule=2pt, rulecolor=\color{codekw},
  xleftmargin=18pt, showstringspaces=false, breaklines=true,
  columns=fullflexible, keepspaces=true, tabsize=4
}
\lstset{style=assignment}

\titleformat{\section}{\Large\sffamily\bfseries\color{head}}{\thesection}{0.6em}{}
\titleformat{\subsection}{\large\sffamily\bfseries\color{head}}{\thesubsection}{0.6em}{}

\pagestyle{fancy}
\fancyhf{}
\fancyhead[L]{\small\sffamily\course}
\fancyhead[R]{\small\sffamily\student}
\fancyfoot[C]{\small\thepage}

\setlength{\parindent}{0pt}
\setlength{\parskip}{0.5em}

% Labelled answer block: \begin{answerbox}{task number} ... \end{answerbox}
\newtcolorbox{answerbox}[1]{enhanced,colback=white,colframe=codekw,boxrule=0.6pt,arc=2pt,
  title={\sffamily\bfseries Answer to Task #1},colbacktitle=codekw,coltitle=white}

\begin{document}

\begin{center}
  {\sffamily\small\color{gray}\course}\\[4pt]
  {\sffamily\LARGE\bfseries\color{head}\assignment}\\[6pt]
  \student \quad\textbar\quad Submitted \today
\end{center}

\begin{tcolorbox}[enhanced,colback=codebg,colframe=codebg,borderline west={3pt}{0pt}{codestr},arc=0pt]
  \textbf{Deliverables.} A single file \texttt{hashtable.py}, this write-up as a PDF, and the output of the test script. Late submissions lose 10 percent per day.
\end{tcolorbox}

\section{Tasks}
% Edit the task list to match your assignment brief.
\begin{enumerate}[label=\textbf{Task \arabic*.},leftmargin=*,align=left]
  \item Implement a hash table with separate chaining that supports \texttt{put}, \texttt{get} and \texttt{remove}.
  \item Resize the table when the load factor exceeds 0.75.
  \item Measure the average time of \texttt{get} for $10^3$ to $10^6$ keys and explain the results.
\end{enumerate}

\section{Implementation}
\subsection{Core class}
\begin{lstlisting}
class HashTable:
    """A hash table using separate chaining."""

    def __init__(self, capacity=8):
        self.buckets = [[] for _ in range(capacity)]
        self.size = 0

    def _index(self, key):
        return hash(key) % len(self.buckets)

    def put(self, key, value):
        bucket = self.buckets[self._index(key)]
        for pair in bucket:
            if pair[0] == key:
                pair[1] = value  # update existing key
                return
        bucket.append([key, value])
        self.size += 1
        if self.size / len(self.buckets) > 0.75:
            self._resize()
\end{lstlisting}

\subsection{Resizing}
Doubling the number of buckets keeps chains short. Every item must be rehashed because the index depends on the table length.
\begin{lstlisting}[firstnumber=20]
    def _resize(self):
        old = self.buckets
        self.buckets = [[] for _ in range(2 * len(old))]
        self.size = 0
        for bucket in old:
            for key, value in bucket:
                self.put(key, value)
\end{lstlisting}

\subsection{Testing}
Run the tests from a terminal with \lstinline[language=bash]!python -m pytest tests/!. Sample output:
\begin{lstlisting}[language={},numbers=none,backgroundcolor=\color{black!85},basicstyle=\ttfamily\small\color{white},rulecolor=\color{black}]
============ 6 passed in 0.04s ============
\end{lstlisting}

\section{Answers}

\begin{answerbox}{1}
All three operations hash the key to a bucket and then scan that bucket. With a good hash function and a bounded load factor the expected chain length is constant, so each operation runs in $O(1)$ expected time.
\end{answerbox}

\begin{answerbox}{2}
The table resizes when \texttt{size / capacity > 0.75}. A single resize costs $O(n)$, but because the capacity doubles each time the amortised cost per insertion remains $O(1)$.
\end{answerbox}

\begin{answerbox}{3}
\begin{center}
\begin{tabular}{rrr}
  \toprule
  Keys & Mean \texttt{get} time (ns) & Resizes \\
  \midrule
  1,000 & 142 & 7 \\
  10,000 & 150 & 11 \\
  100,000 & 161 & 14 \\
  1,000,000 & 188 & 17 \\
  \bottomrule
\end{tabular}
\end{center}
Lookup time stays almost flat, which matches the $O(1)$ expectation. The small increase at one million keys is likely caused by CPU cache misses rather than longer chains.
\end{answerbox}

\section{Reflection}
Write a few sentences about what you found difficult, what you would improve and any sources or help you used.

\end{document}
